\documentclass[11pt,a4paper]{article}
\usepackage[margin=22mm]{geometry}
\usepackage{amsmath,amssymb}
\usepackage{fontspec}
\usepackage{polyglossia}
\setdefaultlanguage{russian}
\setmainfont{DejaVu Serif}
\usepackage[unicode,hidelinks]{hyperref}

\allowdisplaybreaks[1]
\raggedright
\setlength{\parskip}{0.35em}
\newcommand{\dd}{\mathop{}\!\mathrm{d}}
\newcommand{\avg}[1]{\langle #1\rangle}
\newcommand{\ket}[1]{\lvert #1\rangle}
\newcommand{\bra}[1]{\langle #1\rvert}
\newcommand{\braket}[2]{\langle #1\mid #2\rangle}
\DeclareMathOperator{\Var}{Var}

\newcommand{\solutiontitle}[4]{%
  \hypersetup{pdftitle={Квантовая механика 1. Решения ДЗ №#1}}%
  \begin{center}
    {\small Квантовая механика 1 \enspace\textperiodcentered\enspace Семестр 1\par}
    \vspace{1ex}
    {\LARGE\bfseries Решения домашнего задания №#1\par}
    \vspace{0.5ex}
    {\large #2\par}
  \end{center}
  \noindent{\small Условия приведены в файле «#3», #4.}\par\medskip
}

\begin{document}
\solutiontitle{8}{Связанные состояния прямоугольных и дельта-ям}{ДЗ2026(1)\_08 (1).pdf}{с.~1}

\section{Конечная прямоугольная яма}

Пусть \(U_0>0\), \(L=a/2\) и \(y=x-L\). Энергию \(E\) отсчитываем
от потенциала снаружи. При \(E\ge0\) внешние решения осциллируют
или линейны и не дают ненулевого \(L^2\)-решения. Положительность
кинетической энергии и \(U+U_0\ge0\) дают \(E+U_0>0\).
Поэтому для связанного уровня \(-U_0<E<0\).
Введём
\[
 k=\frac{\sqrt{2m(E+U_0)}}{\hbar},\qquad
 \kappa=\frac{\sqrt{-2mE}}{\hbar},\qquad
 z=kL,\quad s=\kappa L,\quad z_0=\frac{L\sqrt{2mU_0}}{\hbar}.
\]
Симметрия позволяет выбрать состояния чётными или нечётными
по \(y\). Внутри они пропорциональны \(f(ky)\), где
\(f=\cos\) или \(f=\sin\). При \(y>L\) остаётся убывающая
экспонента, коэффициент которой фиксирует непрерывность
\(\Psi(y)=A f(z)e^{-\kappa(y-L)}\). Слева берём отражение
с нужной чётностью. Непрерывность производной даёт
\(k f'(z)=-\kappa f(z)\). Значение \(f(z)=0\) потребовало бы
одновременно \(f'(z)=0\), поэтому деление допустимо. Получаем
\[
 s=\sqrt{z_0^2-z^2}=
 \begin{cases}
 z\tan z & \text{чётные состояния},\\
 -z\cot z & \text{нечётные состояния},
 \end{cases}
 \qquad
 E=-\frac{\hbar^2}{2mL^2}\,(z_0^2-z^2).
\]

\textbf{(a)} Для мелкой ямы \(z_0\ll1\) остаётся одна чётная
ветвь. Из \(\kappa L=z\tan z=z^2+O(z^4)\) следует
\(s=O(z_0^2)\) и \(z^2=z_0^2-s^2=z_0^2+O(z_0^4)\).
Значит \(\kappa=(mU_0 a/\hbar^2)[1+O(z_0^2)]\) и
\[
 E=-\frac{mU_0^2a^2}{2\hbar^2}\,[1+O(z_0^2)] .
\]
Положительная энергия из диапазона \(0<E<U_0\) в условии
соответствует отсчёту от дна. При таком отсчёте ответ
\(E_{\text{дно}}=U_0+E\).

\textbf{(b)} При \(U_0\to\infty\) и фиксированном номере
уровня \(n=1,2,\ldots\) величина \(z\) остаётся в интервале
\(((n-1)\pi/2,n\pi/2)\). Левая часть уравнения
с \(\sqrt{z_0^2-z^2}\) растёт без границ, и корень должен
приближаться к правому полюсу. Тогда \(k\to n\pi/a\) и
\[
 E_n=-U_0+\frac{\pi^2\hbar^2 n^2}{2ma^2}+o(1).
\]
Энергия над дном стремится к соответствующему уровню бесконечной ямы.

\textbf{(c)} На каждом интервале
\((j\pi/2,(j+1)\pi/2)\), \(j=0,1,\ldots\), ровно одна из
правых частей положительна. Она возрастает от нуля до бесконечности,
поскольку и \(z\), и выбранная положительная тригонометрическая
функция возрастают. Функция \(\sqrt{z_0^2-z^2}\) убывает.
Если \(j\pi/2<z_0\), пересечение происходит до полюса или до
\(z=z_0\), где левая часть равна нулю. При совпадении \(z_0\)
с левым концом получилось бы
\(E=0\) без нормируемого хвоста. Таким образом ровно \(N\ge1\)
уровней существует при
\[
 \frac{\pi^2\hbar^2(N-1)^2}{2ma^2}<U_0\le
 \frac{\pi^2\hbar^2 N^2}{2ma^2}.
\]

\section{Одиночная дельта-яма}

Пусть \(\alpha>0\), \(E=-\hbar^2\kappa^2/(2m)\), \(\kappa>0\).
По обе стороны от нуля остаются убывающие экспоненты.
Непрерывность даёт \(\Psi=A e^{-\kappa|x|}\).
Интегрирование уравнения Шрёдингера через ноль даёт
\(\Psi'(0^+)-\Psi'(0^-)=-2m\alpha\Psi(0)/\hbar^2\).
Подстановка даёт единственное \(\kappa=m\alpha/\hbar^2\).
Из \(\int e^{-2\kappa|x|}\dd x=1/\kappa\) следуют
\[
 E=-\frac{m\alpha^2}{2\hbar^2},\qquad \Psi(x)=\sqrt\kappa\,e^{-\kappa|x|}.
\]

\textbf{(a)} Пусть \(a(p)\) есть преобразование с ядром
\(e^{-ipx/\hbar}/\sqrt{2\pi\hbar}\). Поскольку
\(\Psi(0)=(2\pi\hbar)^{-1/2}\int a(p')\dd p'\), в импульсах
\[
 \left(\frac{p^2}{2m}+|E|\right)a(p)=
                 \frac{\alpha}{2\pi\hbar}\int a(p')\,\dd p'.
\]
Если правый интеграл ноль, то \(a(p)=0\).
В противном случае делим на скобку, интегрируем и сокращаем
ненулевой интеграл. Так как
\(\int_\mathbb R\dd p/(p^2+\hbar^2\kappa^2)=\pi/(\hbar\kappa)\),
получаем \(1=m\alpha/(\hbar^2\kappa)\).
То же уравнение вместе с нормировкой даёт
\[
 a(p)=\sqrt{\frac{2}{\pi}}\,
              \frac{(\hbar\kappa)^{3/2}}{p^2+(\hbar\kappa)^2}.
\]
Здесь \(\int\dd p/(p^2+b^2)^2=\pi/(2b^3)\) получается
дифференцированием простого интеграла по \(b=\hbar\kappa\).
Обратное преобразование использует
\(\int e^{ipx/\hbar}\dd p/(p^2+b^2)=(\pi/b)e^{-b|x|/\hbar}\)
и даёт найденную \(\Psi\).

\textbf{(b)} Средние \(x,p\) равны нулю по чётности плотностей.
Интегрирование по частям даёт
\(2\kappa\int_0^\infty x^2e^{-2\kappa x}\dd x=1/(2\kappa^2)\).
В импульсном интеграле \(p^2=(p^2+b^2)-b^2\)
даёт \(\int p^2\dd p/(p^2+b^2)^2=\pi/(2b)\).
Итак,
\[
 \Var(x)=\frac{1}{2\kappa^2},\qquad \Var(p)=\hbar^2\kappa^2 .
\]
Произведение стандартных отклонений равно \(\hbar/\sqrt2\).

\section{Яма с одной бесконечной стенкой}

Определим \(k,\kappa\) как в задаче 1, а теперь
\(z=ka,\ s=\kappa a,\ z_0=a\sqrt{2mU_0}/\hbar\).
Стенка требует \(\Psi(0)=0\). Поэтому при \(0<x<a\)
\(\Psi=A\sin kx\), а при \(x>a\) из непрерывности
\(\Psi=A\sin(ka)e^{-\kappa(x-a)}\).
Непрерывность производной при \(a\) даёт \(k\cos(ka)=-\kappa\sin(ka)\).
Нуль \(\sin(ka)\) невозможен одновременно с нулём \(\cos(ka)\).
Получаем спектр
\[
 -z\cot z=\sqrt{z_0^2-z^2},\qquad
 E=-\frac{\hbar^2}{2ma^2}(z_0^2-z^2).
\]
На каждой ветви \(((j+\tfrac12)\pi,(j+1)\pi)\), \(j=0,1,\ldots\),
левая часть растёт от нуля до бесконечности. Правая часть убывает
до нуля при \(z_0\). Поэтому имеется ровно один корень на ветви
при \((j+\tfrac12)\pi<z_0\), и ни одного при обратном неравенстве
или равенстве. В частности, первый уровень возникает при
\(U_0>\pi^2\hbar^2/(8ma^2)\).

\section{Две одинаковые дельта-ямы}

Обозначим \(g=m\alpha/\hbar^2>0,\ a>0\).
При \(E=-\hbar^2\kappa^2/(2m)\) общий вид \(L^2\)-решения с
изломами в \(0,a\) есть линейная комбинация
\(e^{-\kappa|x|},e^{-\kappa|x-a|}\). Действительно, можно выбрать
коэффициенты так, чтобы эти экспоненты воспроизводили скачки
производной. Разность тогда решает \(\Psi''=\kappa^2\Psi\)
на всей оси и в \(L^2\) равна нулю. При отражении \(x\mapsto a-x\)
экспоненты обмениваются, так что для чётности \(s=\pm1\)
относительно середины имеем
\[
 \Psi_s(x)=C_s\bigl(e^{-\kappa_s|x|}+s e^{-\kappa_s|x-a|}\bigr).
\]
В нуле скачок производной равен \(-2\kappa_s C_s\).
Условие скачка \(-2g\Psi_s(0)\) даёт
\[
 \kappa_s=g(1+s e^{-\kappa_s a}),\qquad E_s=-\frac{\hbar^2\kappa_s^2}{2m}.
\]
Условие в \(a\) следует из симметрии. Интеграл перекрытия
экспонент, разбитый на \(x<0,\ 0<x<a,\ x>a\), равен
\(e^{-\kappa_s a}(a+1/\kappa_s)\). Отсюда
\[
 C_s=\left[\frac{2}{\kappa_s}
           \bigl(1+s(1+\kappa_s a)e^{-\kappa_s a}\bigr)\right]^{-1/2}.
\]
Для \(s=+1\) правая часть уравнения убывает от \(2g\) до \(g\),
поэтому положительный корень единственный. Для \(s=-1\) правая
часть \(g(1-e^{-\kappa a})\) строго вогнута, начинается в нуле
с наклоном \(ga\) и стремится к \(g\).
Значит единственный положительный корень существует при \(ga>1\).
При \(ga=1\) остаётся только порог \(\kappa=0\) без связанного уровня.

Для удалённых ям \(ga\gg1\) положим \(\epsilon=e^{-ga}\).
Из уравнения для \(u_s=\kappa_s/g-1\) следует
\(u_s=s\epsilon e^{-ga u_s}\).
В интервале \(|u_s|\le2\epsilon\) правая часть отображает интервал
в себя и имеет производную по модулю \(O(ga\epsilon)<1\).
Разложение экспоненты даёт \(u_s=s\epsilon+O(ga\epsilon^2)\).
Поэтому при \(E_0=-m\alpha^2/(2\hbar^2)\)
\[
 E_s=E_0[1+2s e^{-ga}+O(ga e^{-2ga})],\qquad
 E_--E_+=4|E_0|e^{-ga}+O(|E_0|ga e^{-2ga}).
\]
Чётный уровень ниже нечётного. В пределе оба стремятся
к энергии одной дельта-ямы.

\end{document}
